%------------------------------------------------------------------------------%
\begin{question}
  Encontre e classifique os pontos críticos da função
  \[
    f(x, y) = xy -x^2 - y^2 -2x -2y +4
  \]
  Avalie os valores extremos locais de $f$
\end{question}

%------------------------------------------------------------------------------%
\begin{resolution}{Derivadas parciais}
  \[
    \D{f}{x}
    \pause
    = \D{}{x}\left(xy -x^2 - y^2 -2x -2y +4\right)
    \pause
    = y -2x -2
  \]
  \pause
  \vspace{\baselineskip}
  \[
    \D{f}{y}
    \pause
    = \D{}{y}\left(xy -x^2 - y^2 -2x -2y +4\right)
    \pause
    = x -2y -2
  \]
\end{resolution}

%------------------------------------------------------------------------------%
\begin{resolution}{Pontos críticos}
  As derivadas existem em todos os pontos do plano

  \pause
  \vspace{\baselineskip}
  Buscando os pontos onde plano tangente é horizontal
\end{resolution}

%------------------------------------------------------------------------------%
\begin{resolution}{Resolvendo o sistema}
\begin{columns}
\column{0.5\linewidth}
  \begin{align*}
    \onslide<+->{f_x(x, y) & = y -2x -2 = 0 \\}
    \onslide<+->{f_y(x, y) & = x -2y -2 = 0}
  \end{align*}

  \vspace{0.1\baselineskip}
  \begin{align*}
    \onslide<+->{x -2y -2 & = 0      \\}
    \onslide<+->{x        & = 2y + 2}
  \end{align*}

\column{0.5\linewidth}
  \begin{align*}
    \onslide<+->{y -2x -2         & = 0  \\}
    \onslide<+->{y -2(2y + 2) - 2 & = 0  \\}
    \onslide<+->{y -4y + -6       & = 0  \\}
    \onslide<+->{-3y              & = 6  \\}
    \onslide<+->{y                & = -2}
  \end{align*}

  \[
    \onslide<+->{x = 2y + 2}
    \onslide<+->{= 2(-2) + 2}
    \onslide<+->{= -2}
  \]

  \onslide<+->{Ponto crítico $(-2, -2)$}
\end{columns}
\end{resolution}

%------------------------------------------------------------------------------%
\begin{resolution}{Calculando as derivadas segundas}
  \[
    \pause
    \DS{f}{x}
    \pause
    = \D{}{x}\left(y -2x -2\right)
    \pause
    = -2
  \]
  \[
    \pause
    \DS{f}{y}
    \pause
    = \D{}{y}\left(x -2y -2\right)
    \pause
    = -2
  \]
  \[
    \pause
    \DM{f}{x}{y}
    \pause
    = \D{}{x}\left(x -2y -2\right)
    \pause
    = 1
  \]

  \pause
  Discriminante
  \[
    \pause
    D(x, y)
    \pause
    = f_{xx}(x, y) f_{yy}(x, y) - f^2_{xy}(x, y)
    \pause
    = (-2)(-2) - 1
    \pause
    = 3
  \]
  \[
  \pause
    D(-2, -2)
    \pause
    = 3
  \]
\end{resolution}

%------------------------------------------------------------------------------%
\begin{resolution}{Caracterização do ponto}
  Como \(D(-2, -2)=3>0\)
  \pause
  e  \(f_{xx}(-2, -2) = -2 < 0\)

  \pause
  \vspace{\baselineskip}
  O ponto \((-2, -2)\) é um ponto de máximo local

  \pause
  \vspace{\baselineskip}
  O valor da função neste máximo local é
  \[
    f(-2, -2)
    \pause
    = \left(xy -x^2 - y^2 -2x -2y +4\right)\evalat{(-2, -2)}{}
  \]
  \[
    \hphantom{f(-2, -2)}
    \pause
    = (-2)(-2) - (-2)^2 - (-2)^2 -2(-2) -2(-2) + 4
  \]
  \[
    \hphantom{f(-2, -2)}
    \pause
    = 4 -4 -4 +4 +4 + 4
  \]
  \[
    \hphantom{f(-2, -2)}
    \pause
    = 8
  \]
\end{resolution}

%------------------------------------------------------------------------------%
